\rtmtight
\begin{tblr}{width=\textwidth,colspec={|Q[c,m,2.1cm]|X[l,m]|},hlines,vlines,hspan=minimal,rowsep=1.5pt,colsep=3pt,row{1}={bg=headgray,font=\bfseries}}
\SetCell{c,m}\hfil\logo\hfil & \textbf{POLITEKNIK NEGERI MALANG}\par\textbf{JURUSAN TEKNOLOGI INFORMASI}\par\textbf{PROGRAM STUDI: D4 TEKNIK INFORMATIKA} \\
\SetCell[c=2]{l} \textbf{RENCANA TUGAS MAHASISWA (RTM) DAN RUBRIK PENILAIAN} & \\
Nama Mata Kuliah & Matematika Dasar \\
Kode Mata Kuliah & RTI251004 \quad \textbf{sks / Jam perminggu:} 2 SKS/4 jam \quad \textbf{Semester:} 1 \\
Dosen Pengampu & \begin{enumerate}\item Muhammad Afif Hendrawan, S.Kom., M.T.\item Prof. Dr. Eng. Cahya Rahmad, S.T., M.Kom.\item Habibie Ed Dien, S.Kom., M.T.\item Erfan Rohadi, S.T., M.Eng., Ph.D.\item Vipkas Al Hadid Firdaus, S.T., M.T.\end{enumerate} \\
\end{tblr}
\assessmentblock{Tugas Latihan Soal Matematika Dasar}{Tugas latihan soal mingguan}{SCPMK0902-00401-SCPMK0902-00404}{Mahasiswa mengerjakan latihan soal mingguan tentang logika, himpunan, relasi dan fungsi, sistem bilangan, induksi matematika, aljabar Boolean, kombinatorial, graf, dan pohon.}{Mengerjakan latihan soal terstruktur di kelas dan take-home, menuliskan langkah penyelesaian secara runtut, membahas jawaban di papan tulis, dan mengumpulkan berkas melalui LMS.}{Lembar jawaban/berkas tugas dengan langkah penyelesaian lengkap yang dikumpulkan melalui LMS.}{Ketepatan prosedur penyelesaian soal logika dan himpunan & 5 & Tabel kebenaran, hukum logika, dan operasi himpunan digunakan dengan langkah runtut dan jawaban akhir benar. & Prosedur umumnya benar, tetapi terdapat kesalahan hitung atau langkah yang terlewat pada sebagian soal. & Prosedur yang dipilih keliru atau jawaban ditulis tanpa langkah penyelesaian. \\ \\
Ketepatan perhitungan relasi, fungsi, dan konversi bilangan & 4.5 & Sifat relasi/fungsi ditentukan benar dan konversi biner, oktal, desimal, heksadesimal dihitung tanpa kesalahan. & Sebagian besar perhitungan benar, tetapi terdapat kesalahan minor pada konversi atau penentuan sifat relasi/fungsi. & Perhitungan banyak yang keliru atau aturan konversi antarbasis tidak digunakan dengan benar. \\ \\
Ketepatan penerapan induksi, aljabar Boolean, dan kombinatorial & 6 & Pembuktian induksi lengkap (basis dan langkah induksi), fungsi Boolean disederhanakan benar, dan rumus kombinatorial dipilih serta dihitung tepat. & Konsep dipilih benar, tetapi terdapat langkah pembuktian atau perhitungan yang belum lengkap. & Pembuktian atau perhitungan salah konsep, tidak selesai, atau hasil akhir tidak dapat ditelusuri. \\ \\
Ketepatan penyelesaian persoalan graf dan pohon & 4.5 & Representasi graf/pohon digambar benar serta lintasan, sirkuit, atau spanning tree ditentukan dengan tepat. & Representasi benar, tetapi penentuan lintasan, sirkuit, atau pohon masih memuat kesalahan. & Representasi keliru atau solusi tidak menjawab persoalan yang diberikan. \\}

\assessmentblock{Kuis 1 Logika dan Himpunan}{Kuis (ujian tulis)}{SCPMK0902-00401}{Mahasiswa mengerjakan kuis individual yang mencakup materi pertemuan 1 s.d. 3: logika proposisi, hukum logika, dan teori himpunan.}{Mengerjakan soal kuis secara individual dalam waktu terbatas dengan menuliskan langkah penyelesaian dan jawaban akhir.}{Lembar jawaban kuis dengan langkah penyelesaian dan jawaban akhir.}{Ketepatan penerapan hukum logika proposisi & 3 & Nilai kebenaran, ekivalensi, tautologi/kontradiksi ditentukan benar dengan hukum logika dan De Morgan yang diterapkan tepat. & Sebagian besar jawaban benar, tetapi terdapat kesalahan penerapan hukum pada satu-dua soal. & Hukum logika diterapkan keliru atau jawaban tidak konsisten dengan tabel kebenaran. \\ \\
Ketepatan operasi himpunan dan diagram Venn & 3 & Operasi himpunan, kardinalitas, dan diagram Venn dikerjakan benar termasuk prinsip inklusi-eksklusi. & Operasi dasar benar, tetapi terdapat kesalahan pada inklusi-eksklusi atau pembacaan diagram. & Operasi himpunan banyak keliru atau diagram tidak merepresentasikan soal. \\ \\
Kejelasan langkah dan interpretasi hasil & 1.5 & Langkah penyelesaian runtut, notasi benar, dan hasil diinterpretasikan sesuai konteks soal. & Langkah dapat diikuti, tetapi notasi atau interpretasi hasil kurang tepat di sebagian jawaban. & Jawaban tanpa langkah yang dapat ditelusuri atau interpretasi hasil keliru. \\}

\assessmentblock{UTS Logika, Himpunan, Relasi, Fungsi, dan Sistem Bilangan}{UTS (ujian tulis)}{SCPMK0902-00401, SCPMK0902-00402}{Mahasiswa mengerjakan ujian tengah semester individual yang mencakup materi pertemuan 1 s.d. 7.}{Mengerjakan soal UTS secara individual dalam waktu terbatas dengan menuliskan langkah penyelesaian yang runtut untuk setiap soal.}{Lembar jawaban UTS dengan langkah penyelesaian dan jawaban akhir.}{Ketepatan penyelesaian soal logika proposisi & 8 & Tabel kebenaran, penyederhanaan proposisi, dan pembuktian ekivalensi dikerjakan benar dan lengkap. & Prosedur benar, tetapi terdapat kesalahan hitung atau langkah penyederhanaan yang terlewat. & Prosedur keliru atau kesimpulan tidak sesuai hasil tabel kebenaran. \\ \\
Ketepatan operasi dan pembuktian himpunan & 7 & Operasi himpunan, hukum himpunan, dan prinsip inklusi-eksklusi diterapkan benar pada seluruh soal. & Sebagian besar operasi benar, tetapi pembuktian atau penerapan inklusi-eksklusi belum lengkap. & Operasi dan pembuktian himpunan banyak yang keliru atau tidak dikerjakan. \\ \\
Ketepatan perhitungan relasi dan fungsi & 8 & Jenis relasi, nilai fungsi, fungsi rekursif, dan kardinalitas ditentukan benar dengan langkah jelas. & Sebagian besar jawaban benar, tetapi terdapat kesalahan pada fungsi rekursif atau penentuan sifat relasi. & Konsep relasi/fungsi tertukar atau perhitungan tidak menghasilkan jawaban benar. \\ \\
Ketepatan konversi dan operasi sistem bilangan & 7 & Konversi biner, oktal, desimal, heksadesimal serta operasi antarbasis dihitung benar seluruhnya. & Konversi umumnya benar, tetapi terdapat kesalahan minor pada operasi antarbasis. & Aturan konversi tidak dikuasai sehingga sebagian besar jawaban salah. \\}

\assessmentblock{Kuis 2 Induksi, Aljabar Boolean, dan Kombinatorial}{Kuis (ujian tulis)}{SCPMK0902-00403}{Mahasiswa mengerjakan kuis individual yang mencakup materi pertemuan 9 s.d. 12: induksi matematika, aljabar Boolean, dan kombinatorial.}{Mengerjakan soal kuis secara individual dalam waktu terbatas dengan menuliskan langkah penyelesaian dan jawaban akhir.}{Lembar jawaban kuis dengan langkah penyelesaian dan jawaban akhir.}{Ketepatan pembuktian dengan induksi matematika & 2.5 & Basis induksi, hipotesis, dan langkah induksi dituliskan lengkap sehingga pembuktian sahih. & Struktur pembuktian benar, tetapi langkah aljabar pada tahap induksi belum lengkap. & Pembuktian melompat, tidak memuat basis/langkah induksi, atau salah konsep. \\ \\
Ketepatan penyederhanaan fungsi Boolean & 2.5 & Fungsi Boolean disusun dari tabel dan disederhanakan benar menggunakan teorema dasar. & Fungsi tersusun benar, tetapi penyederhanaan belum minimal atau ada kesalahan kecil. & Fungsi tidak sesuai tabel atau penyederhanaan melanggar teorema. \\ \\
Ketepatan perhitungan permutasi dan kombinasi & 2.5 & Kaidah pencacahan, permutasi, dan kombinasi dipilih sesuai kasus dan dihitung benar. & Rumus dipilih benar, tetapi terdapat kesalahan hitung pada sebagian soal. & Rumus tertukar antara permutasi dan kombinasi atau hasil hitung salah. \\}

\assessmentblock{UAS Matematika Diskrit}{UAS (ujian tulis)}{SCPMK0902-00403, SCPMK0902-00404}{Mahasiswa mengerjakan ujian akhir semester individual yang mencakup materi pertemuan 1 s.d. 16 dengan penekanan pada induksi, aljabar Boolean, kombinatorial, graf, dan pohon.}{Mengerjakan soal UAS secara individual dalam waktu terbatas dengan menuliskan langkah penyelesaian yang runtut untuk setiap soal.}{Lembar jawaban UAS dengan langkah penyelesaian dan jawaban akhir.}{Ketepatan penerapan induksi dan aljabar Boolean & 9 & Pembuktian induksi sahih dan lengkap serta fungsi Boolean disederhanakan benar sesuai teorema dasar. & Struktur pembuktian dan penyederhanaan benar, tetapi terdapat langkah yang belum lengkap. & Pembuktian tidak sahih atau penyederhanaan Boolean salah konsep. \\ \\
Ketepatan perhitungan kombinatorial & 8.5 & Kaidah pencacahan, permutasi, dan kombinasi (termasuk perulangan objek) dipilih dan dihitung benar pada semua kasus. & Rumus dipilih benar, tetapi terdapat kesalahan hitung atau kasus perulangan yang terlewat. & Pemilihan rumus keliru sehingga hasil pencacahan tidak benar. \\ \\
Ketepatan representasi dan analisis graf & 9 & Graf direpresentasikan benar serta lintasan/sirkuit Euler dan Hamilton ditentukan dengan justifikasi tepat. & Representasi benar, tetapi penentuan lintasan/sirkuit belum disertai justifikasi lengkap. & Representasi graf keliru atau lintasan/sirkuit yang diberikan tidak valid. \\ \\
Ketepatan penyelesaian persoalan pohon & 8.5 & Spanning tree, pohon berakar, dan pohon biner dibangun serta dianalisis benar sesuai definisi. & Konstruksi pohon umumnya benar, tetapi analisis sifat atau traversal masih ada kesalahan. & Konstruksi pohon tidak sesuai definisi atau tidak menjawab persoalan. \\}

\begin{tblr}{width=\textwidth,colspec={|Q[l,m,3.2cm]|X[l,m]|},hlines,vlines,hspan=minimal,row{1-3}={bg=headgray},rowsep=1.5pt,colsep=3pt}
Jadwal Pelaksanaan: & Minggu 1--17 sesuai RPS. Setiap evaluasi memiliki RTM dan rubrik multi-kriteria. \\
Lain-Lain Yang Diperlukan: & LMS, lembar soal dan jawaban, perangkat lunak sesuai mata kuliah, dan perangkat presentasi. \\
Pustaka: & Mengacu pustaka utama dan pendukung pada bagian RPS. \\
\end{tblr}
